\documentclass{beamer}
\usetheme{Madrid}

% Encoding and fonts for pdfLaTeX
\usepackage[utf8]{inputenc}
\usepackage[T1]{fontenc}
\usepackage{lmodern}
\usepackage{hyperref}

\title{Prompt Engineering}
\author{Mehmet Ali Akyol, PhD}
\institute{Ankara Medipol University}
\date{\today}

\begin{document}

\begin{frame}
    \titlepage
    \begin{center}
        \vspace{1cm}
        \textbf{Week 5: Controlling Style, Tone, and Format}
    \end{center}
\end{frame}

\section{Overview}

\begin{frame}
    \frametitle{Where We Are}
    \begin{itemize}
        \item Prior weeks established fundamentals and reusable pattern frameworks.
        \item \textbf{This week:} Focus on precise control of style, tone, register, and structured formatting without rehashing basics.
    \end{itemize}
\end{frame}

\begin{frame}
    \frametitle{Weekly Learning Objectives}
    \begin{itemize}
        \item Rapidly profile audience and medium to derive tone/register parameters.
        \item Operationalize voice guidelines (brand, compliance, cultural sensitivity) in compact prompt clauses.
        \item Specify parse-ready output (tables / JSON / labeled sections) aligned to downstream pipeline needs.
        \item Instrument prompts for self-audit to catch tone, register, or formatting drift.
        \item Produce lightweight, reusable micro style guides for collaborative human–AI drafting.
    \end{itemize}
\end{frame}

\begin{frame}
    \frametitle{Key Concepts}
    \begin{itemize}
        \item Tone ladder calibration vs.\ tone adjectives
        \item Voice fingerprint (syntax rhythm, perspective, lexical field)
        \item Register levers (jargon density, sentence length, metaphor frequency)
        \item Structural contracts (schema / heading spec / placeholder tokens)
        \item Drift detection and corrective refinement loops
    \end{itemize}
\end{frame}

\section{Controlling Style and Tone}


\begin{frame}[fragile]
    \frametitle{Quick Question: Which Email Would You Send?}
    \begin{columns}[T]
        \column{0.48\textwidth}
        \textbf{Version A}
        \begin{block}{}
        {\small
        Subject: Payment Due
        
        Your invoice is overdue. Please pay as soon as possible. If you have questions, contact us.
        
        Thanks.
        }
        \end{block}
        
        \column{0.48\textwidth}
        \textbf{Version B}
        \begin{block}{}
        {\small
        Subject: Friendly Payment Reminder
        
        Hi [Name],
        
        We haven't received payment for Invoice \#2847 (due Nov 1). We understand things get busy!
        
        Could you please submit payment by Nov 10? Reply if you need help.
        
        Thank you for your business!
        }
        \end{block}
    \end{columns}
    \vspace{0.3cm}
    \textbf{Quick poll:} A, B, or Depends?
\end{frame}

\begin{frame}
    \frametitle{Linguistic Register: Beyond Formal vs. Casual}
    \textbf{Three dimensions of register control:}
    \begin{itemize}
        \item \textbf{Field (Topic Domain):} Controls terminology density and specialization.
        \begin{itemize}
            \item Technical: \emph{myocardial infarction, hemodynamic instability}
            \item Lay: \emph{heart attack, blood pressure problems}
        \end{itemize}
        \item \textbf{Tenor (Relationship):} Manages power distance and social connection.
        \begin{itemize}
            \item Authority: \emph{You must submit\ldots}, \emph{I recommend\ldots}
            \item Peer: \emph{Let's review\ldots}, \emph{What do you think?}
            \item Subordinate: \emph{Would it be possible\ldots}, \emph{I'd appreciate\ldots}
        \end{itemize}
        \item \textbf{Mode (Channel):} Adjusts for delivery style and planning.
        \begin{itemize}
            \item Written-planned: complete sentences, cohesive structure
            \item Spoken-spontaneous: fragments, fillers, self-correction
        \end{itemize}
    \end{itemize}
\end{frame}

\begin{frame}
    \frametitle{Multi-Dimensional Register Specifications}
    \textbf{Instead of vague instructions:}
    \begin{itemize}
        \item ``Be professional'' $\rightarrow$ ambiguous
        \item ``Make it friendly'' $\rightarrow$ inconsistent results
    \end{itemize}
    
    \vspace{0.3cm}
    \textbf{Use precise dimensional specs:}
    \begin{itemize}
        \item \texttt{CEFR B1 + peer tenor + written-planned mode}
        \item \texttt{Domain: healthcare lay + subordinate tenor + email mode}
        \item \texttt{Technical depth: expert + authority tenor + monologic}
    \end{itemize}
    
    \vspace{0.3cm}
    \textbf{Result:} Precise, reproducible style control without trial-and-error.
\end{frame}

\begin{frame}
    \frametitle{Voice and Tone Diagnostics}
    \textbf{Step 1: Audience Analysis}
    \begin{itemize}
        \item Industry context: legal (risk-averse), tech (innovation-forward), healthcare (empathetic)
        \item Reading level: CEFR B1 (general public) vs.\ C2 (domain experts)
        \item Cultural norms: directness (US/Germany) vs.\ indirectness (Japan/Middle East)
    \end{itemize}
    
    \textbf{Step 2: Tone Mapping}
    \begin{itemize}
        \item Vague adjectives $\rightarrow$ Observable behaviors
        \item \emph{Warm}: Use names, express gratitude, acknowledge effort
        \item \emph{Urgent}: Short sentences, action verbs, explicit deadlines
        \item \emph{Neutral}: Third-person, passive constructions, hedge words
    \end{itemize}
\end{frame}

\begin{frame}
    \frametitle{Voice and Tone Diagnostics (continued)}
    \textbf{Step 3: No-Go Zones}
    \begin{itemize}
        \item Banned words: \emph{cheap, guarantee, revolutionary, disrupt}
        \item Prohibited claims: \emph{We're the best, FDA-approved (if not true)}
        \item Brand sensitivities: competitor names, controversial topics
    \end{itemize}
    
    \vspace{0.3cm}
    \textbf{Step 4: Provide Exemplars}
    \begin{itemize}
        \item \textbf{Strong sample:} \emph{Thanks for your patience. We've resolved the issue and added extra support hours.}
        \item \textbf{Counter-example:} \emph{Your ticket was processed. Issue closed.}
        \item \textbf{Why it matters:} AI learns boundaries from contrasts, not just positive examples.
    \end{itemize}
\end{frame}

\begin{frame}
    \frametitle{Brand Voice Fingerprinting}
    \begin{block}{Concept}
        Extract measurable, reproducible voice attributes from existing content to ensure AI outputs match your brand.
    \end{block}
    
    \textbf{Four dimensions to capture:}
    \begin{enumerate}
        \item \textbf{Syntax signature:} Measure sentence structure patterns
        \item \textbf{Lexical field:} Identify preferred vocabulary choices
        \item \textbf{Rhythm:} Analyze pacing and variety
        \item \textbf{Perspective:} Determine point-of-view consistency
    \end{enumerate}
\end{frame}

\begin{frame}
    \frametitle{Brand Voice Fingerprinting: Measurable Attributes}
    \begin{columns}[T]
        \begin{column}{0.48\textwidth}
            \textbf{Syntax Signature}
            \begin{itemize}
                \item Avg.\ sentence length: 12--18 words
                \item Clause complexity: 1.5 clauses/sentence
                \item Passive voice: $<$10\%
            \end{itemize}
            
            \vspace{0.2cm}
            \textbf{Lexical Field}
            \begin{itemize}
                \item Preferred: \emph{enable, drive, foster}
                \item Avoid: \emph{leverage, synergy, utilize}
            \end{itemize}
        \end{column}
        
        \begin{column}{0.48\textwidth}
            \textbf{Rhythm}
            \begin{itemize}
                \item Sentence length variance: high
                \item Fragment usage: occasional for emphasis
                \item Lists: 3--5 items, parallel structure
            \end{itemize}
            
            \vspace{0.2cm}
            \textbf{Perspective}
            \begin{itemize}
                \item First-person plural: ``We help you\ldots''
                \item vs.\ Institutional: ``The platform provides\ldots''
            \end{itemize}
        \end{column}
    \end{columns}
    
    \vspace{0.3cm}
    \textbf{Process:} Analyze 3--5 approved samples $\rightarrow$ Extract metrics $\rightarrow$ Encode as prompt constraints
\end{frame}

\begin{frame}
    \frametitle{Style Guides in Prompts}
    \begin{itemize}
        \item Embed condensed brand voice manuals: values, vocabulary, cadence.
        \item Use bullet checklists to anchor style: ``Always include\ldots'', ``Never include\ldots''.
    \item Request self-audit — \emph{List 3 style guidelines you followed before the final answer.} 
        \item Maintain iterative memory with conversation tags (e.g., \texttt{StyleGuide v2}).
    \end{itemize}
\end{frame}

\begin{frame}
    \frametitle{Tone Shift Mechanics: Instruction Levers}
    \begin{block}{Core Insight}
        Small, atomic prompt changes yield predictable tone movement.
    \end{block}
    \begin{itemize}
        \item \textbf{Contraction control:} \emph{Avoid contractions} \textrightarrow{} formal; \emph{Use contractions} \textrightarrow{} casual.
        \item \textbf{Sentence directive:} \emph{Use short, direct sentences} vs.\ \emph{Use complex, flowing prose}.
        \item \textbf{Affective markers:} \emph{Include gratitude and empathy} vs.\ \emph{Use neutral, objective language}.
        \item \textbf{Urgency cues:} \emph{Emphasize deadline and consequences} vs.\ \emph{Provide context without pressure}.
        \item Test each lever in isolation to measure effect size.
    \end{itemize}
\end{frame}

\begin{frame}
    \frametitle{Register and Reading Level}
    \begin{itemize}
        \item Specify reading grade (CEFR level, Flesch-Kincaid, Common European levels).
        \item Ask for analogies or metaphors suited to audience background.
        \item Include domain-specific vocabulary lists to use or avoid.
    \item Prompt for translanguaging — \emph{Introduce one Turkish idiom and translate it.} 
    \end{itemize}
\end{frame}

\section{Formatting Mastery}

\begin{frame}
    \frametitle{Structural Contracts: Parse-Ready Output}
    \begin{block}{Concept}
        A structural contract is an explicit, testable specification for output shape.
    \end{block}
    \begin{itemize}
        \item \textbf{Why?} Downstream automation (APIs, DBs, dashboards) requires predictable format.
        \item \textbf{Components:} section labels, field names, data types, ordering rules.
        \item \textbf{Enforcement:} request model to echo the contract before generating content.
        \item \textbf{Validation hook:} \emph{If any required field is missing, output INCOMPLETE and list missing fields}.
    \end{itemize}
\end{frame}

\begin{frame}
    \frametitle{Structured Output Techniques}
    \begin{itemize}
        \item Define exact headings, table columns, or JSON fields.
        \item Use placeholders: \texttt{<Intro>}, \texttt{<Key Risks>}, \texttt{<Call to Action>}.
        \item Enforce ordering: numbered lists, priority ranking, timeline progression.
    \item Incorporate validation prompts — \emph{Confirm all fields are filled; note any missing data.} 
    \end{itemize}
\end{frame}

\begin{frame}[fragile]
    \frametitle{Schema Enforcement}
    \begin{itemize}
        \item Provide schema with types and constraints.
        \item Ask model to restate schema pre-output to ensure alignment.
        \item Encourage graceful failure — \emph{If you cannot comply, return VALIDATION\_ERROR}.
    \end{itemize}
\vspace{0.2cm}
\textbf{Example}
\begin{verbatim}
Return JSON array of objects with keys:
- task (string, sentence case)
- tone (one of: neutral, upbeat, urgent)
- word_limit (integer, 50-200)
- notes (array of 1-3 short strings)
Only output JSON.
\end{verbatim}
\end{frame}

\begin{frame}
    \frametitle{Multi-Channel Formatting}
    \begin{itemize}
        \item Build variants: email, SMS, chatbot, longform doc.
    \item Use conditional instructions — \emph{If channel=SMS, limit to 160 characters.} 
        \item Include metadata for automation: tags, CTA keywords, compliance flags.
        \item Request summary + main content for editorial review workflows.
    \end{itemize}
\end{frame}

\section{Quality Control}

\begin{frame}
    \frametitle{Multi-Stakeholder Voice Reconciliation}
    \begin{block}{Challenge}
        Different stakeholders (legal, marketing, product) impose conflicting voice rules.
    \end{block}
    \begin{itemize}
        \item \textbf{Strategy 1 — Layered prompts:} Legal disclaimers as append-only; marketing tone for body.
        \item \textbf{Strategy 2 — Priority matrix:} Encode rule precedence (compliance > brand > preference).
        \item \textbf{Strategy 3 — Modular outputs:} Generate compliant core + stylized wrapper separately; merge post-hoc.
        \item Document conflicts in a decision log; iterate with stakeholder input.
    \end{itemize}
\end{frame}

\begin{frame}
    \frametitle{Detecting Drift}
    \begin{itemize}
        \item Prompt for reflection — \emph{Describe how this output aligns with tone X}.
        \item Use comparison prompts — \emph{Compare to sample and note deviations}.
        \item Run batch consistency checks with templated critique prompts.
        \item Employ versioning: keep changelog of style guide tweaks.
    \end{itemize}
\end{frame}

\begin{frame}
    \frametitle{Revision Loops}
    \begin{itemize}
        \item Phase 1: Draft with full instructions.
        \item Phase 2: Critique focusing on tone, compliance, formatting.
        \item Phase 3: Revise using critique as structured feedback.
        \item Escalate to human review when checks fail twice.
    \end{itemize}
\end{frame}

\begin{frame}
    \frametitle{Programmatic Tone Testing}
    \begin{block}{Concept}
        Automate tone evaluation by embedding test prompts in your pipeline.
    \end{block}
    \begin{itemize}
        \item \textbf{Self-labeling:} Ask model to classify its own tone (formal / neutral / casual).
        \item \textbf{Reference comparison:} Provide gold-standard sample; request deviation score (0--10).
        \item \textbf{Sentiment analysis:} Use lightweight NLP tool to measure valence and arousal.
        \item \textbf{Readability metrics:} Flesch-Kincaid, SMOG index for grade-level verification.
        \item Run tests on batch outputs; flag outliers for manual review.
    \end{itemize}
\end{frame}

\section{Worked Examples}

\begin{frame}
    \frametitle{Healthcare: Patient Outreach}
    \begin{itemize}
        \item Audience: adult patients, B1 reading level.
        \item Tone: empathetic, reassuring; include disclaimer.
        \item Format: 3 bullets + call to action button text.
        \item Validation: highlight any missing appointment data.
    \end{itemize}
    \vspace{0.3cm}
    \textbf{Before:} Generic, no tone control, missing validation.\\
    \textbf{After:} Empathy markers, B1 vocabulary, graceful failure for missing fields.
\end{frame}

\begin{frame}
    \frametitle{Finance: Investor Update}
    \begin{itemize}
        \item Voice: confident, transparent, data-led.
        \item Include table (metric, Q2 value, delta vs.\ Q1).
        \item Provide macro context paragraph and risk disclosure statement.
        \item Request internal summary for leadership in 100 words.
    \end{itemize}
    \vspace{0.3cm}
    \textbf{Before:} Unstructured narrative, no quantitative anchor.\\
    \textbf{After:} Multi-format output (table + paragraphs), dual audiences (investors + internal).
\end{frame}

\begin{frame}
    \frametitle{Education: Lesson Plan Variants}
    \begin{itemize}
        \item Generate two versions: 6th grade science, undergraduate seminar.
        \item Tone shift from playful to scholarly.
        \item Format: table with sections (Hook, Core Content, Activity, Assessment).
        \item Reflection: note adjustments made for each audience.
    \end{itemize}
    \vspace{0.3cm}
    \textbf{Before:} Single-register output, no audience differentiation.\\
    \textbf{After:} Parallel outputs with explicit register shift + meta-commentary on changes.
\end{frame}

\begin{frame}
    \frametitle{Legal-Marketing Hybrid: Product Announcement}
    \begin{block}{Scenario}
        Legal requires risk disclaimers; marketing wants upbeat, benefit-forward language.
    \end{block}
    \begin{itemize}
        \item \textbf{Approach:} Two-pass generation.
        \item Pass 1: Marketing-optimized body (tone: upbeat; focus: user benefits).
        \item Pass 2: Append legal footer (tone: neutral; format: separate section with \texttt{<Legal>} tag).
        \item \textbf{Validation:} Confirm legal section is non-empty and placed after marketing content.
    \end{itemize}
    \vspace{0.2cm}
    \textbf{Outcome:} Satisfies both stakeholders without tone bleed between sections.
\end{frame}

\section{Activities}

\begin{frame}
    \frametitle{Activity 1: Tone Ladder Workshop}
    \begin{block}{Goal}
        Practice moving a message up and down the formality scale.
    \end{block}
    \begin{itemize}
        \item Choose a baseline message (e.g., payment reminder).
        \item Produce five variants from \emph{stern} to \emph{friendly} using prompt tweaks.
        \item Annotate which instructions caused each shift.
    \end{itemize}
\end{frame}

% (Activity 2 removed per course adjustment)
% (Assignment and rubric removed to avoid redundancy with earlier weeks)

\section{Wrap-up}

\begin{frame}
    \frametitle{Key Takeaways}
    \begin{itemize}
        \item Tone and format control start with precise diagnostics.
        \item Combine style guides with structural constraints for reliability.
        \item Revision prompts catch drift and safeguard brand voice.
        \item Documented playbooks enable scalable human-AI collaboration.
    \end{itemize}
\end{frame}

\begin{frame}{Questions and Discussion}
    \begin{center}
    {\Huge Thank You!}
    \vspace{1cm}

        \textbf{Dr.\ Mehmet Ali Akyol}\\
        \texttt{mehmetali.akyol@gmail.com}
    \vspace{1cm}

    {\large Questions?}
    \end{center}
\end{frame}

\end{document}
